

\section{Related work}\label{sec:background_and_related_work}

\xhdr{Autoregressive structured prediction in NLP}
%Structured prediction refers to the process of solving inference tasks that involve producing structured outputs.
%For instance, the output structure in cIE is a set of subject-relation-object triplets separated by some pre-defined special tokens.
%The output structure in CP is a linearized parse tree.
%
It has become popular to use autoregressive generative models for structured prediction tasks, as this fits the training mode and specific strengths of language models~\cite{NIPS2015_grammar, athiwaratkun-etal-2020-augmented, liu-etal-2022-autoregressive, Paolini2021}.
%
For instance, \citet{NIPS2015_grammar} modeled dependency parsing as a sequence generation problem and leveraged LSTMs to tackle it effectively.
\citet{decao2021autoregressive} proposed autoregressive language models to address entity linking, entity disambiguation, and document retrieval tasks.

\xhdr{Constrained decoding}
For tasks where the output needs to satisfy certain constraints, constrained decoding has been proposed to guide the generation process to produce valid outputs.
%
For instance, \citet{hokamp-liu-2017-lexically, hu-etal-2019-improved, post-vilar-2018-fast} proposed lexically\hyp constrained sequence decoding for generation tasks.
%
\citet{anderson-etal-2017-guided} extended the beam search algorithm by pruning the search space based on constraints.
%
\citet{scholak-etal-2021-picard} leveraged an incremental parsing technique to generate approximately valid SQL queries.
%\cite{liu2022autoregressive_ETHZ} suggested an action-based approach to generate valid structured outputs.
%
\citet{decao2021autoregressive} addressed entity disambiguation with trie-based lexical  constraints at decoding time to force outputs to be valid entities.
\citet{josifoski-etal-2022-genie} 
addressed closed information extraction by combining trie-based lexical constraints with state-based constraints to force the output to be valid triplet sequences.
\citet{bailin_wang2023grammar} used Earley parser--based constraints to force outputs to lie in a domain-specific language.

\xhdr{Grammar-constrained decoding}
\citet{deutsch-etal-2019-general} proposed a general framework for push-down automata--based constraints and applied it to parsing tasks, which is equivalent to CFG-based constraints in terms of expressiveness.
\citet{yin-neubig-2017-syntactic} proposed a grammar-powered neural architecture for general-purpose code generation.
\citet{shin-etal-2021-constrained} proposed to use grammar-constraints to solve semantic parsing tasks with GPT-3 and envisioned that a semantic parser can be built by combining a large language model with a carefully designed grammar.
\citet{roy2022benchclamp} released a toolkit for grammar-constrained decoding to generate valid meaning representations.
Finally, \citet{stengeleskin2023zero} tested the ability of LLMs to solve parsing task under ambiguity and used grammar constraints to ensure the grammaticality of the output.
% insert a section called method


%While the output structure in tasks like ED and cIE is simpler than in parsing or code generation, there are two additional challenges in these tasks:
%(1) the constraints could be input-dependent, which means that the constraints are different for different inputs;
%(2) the constraints have millions of rules from the knowledge base, posing a scalability challenge.

